% ID: SEL-D0030-OB-E04
% ESTADO: cotejada. Segunda pasada visual completa y borrador compilado revisados.
% Fuente: Selectividad/matematicas_aplicadas/titular_junio/Examen-Matemáticas Aplicadas a las Ciencias Sociales.pdf
% SHA-256: 5bd7838829b9f653746319907f86642017d12e66ed097db0126fb544a2d0f632
% Página: 2; opción B; ejercicio original 4.

En un estudio sobre la utilización de nuevas tecnologías entre los estudiantes de Bachillerato, se ha realizado una encuesta a 500 estudiantes elegidos mediante muestreo aleatorio simple, resultando que 380 de ellos son usuarios de una determinada red social.
\begin{enumerate}[label=\alph*)]
\item \textbf{(1.5 puntos)} Calcule un intervalo de confianza al 97\% para la proporción de estudiantes que son usuarios de esa red social.
\item \textbf{(1 punto)} Suponiendo que se mantiene la proporción muestral, determine el número mínimo de estudiantes a los que sería preciso entrevistar para que, con un nivel de confianza del 96\%, el error cometido al estimar la proporción de usuarios de la citada red social no supere el 2\%.
\end{enumerate}
